\chapter{Introdução}
\label{cap:introducao}
	
Os Grandes Modelos de Linguagem (do inglês \textit{Large Language Models}, LLMs) têm se consolidado como solucionadores de propósito geral para tarefas de linguagem natural, impulsionando sua adoção em sistemas reais em produção, desde assistentes conversacionais até ferramentas de apoio à decisão em ambientes institucionais \cite{zhao2023survey}. No setor público brasileiro, esse movimento é especialmente visível: o Poder Judiciário já opera sistemas baseados em LLMs para auxiliar na análise de documentos processuais \cite{trt4galileu2026}. Contudo, a adoção em larga escala desses modelos em contextos sensíveis traz consigo uma exigência regulatória e técnica que não pode ser ignorada. A Resolução CNJ n\textordmasculine{} 615/2025 veda o uso de LLMs externos para o processamento de dados sigilosos e determina que só devem ser adotadas soluções de código aberto ou comerciais que permitam flexibilidade de adaptação aos contextos locais \cite{cnj615}, o que na prática significa que tribunais e órgãos públicos precisam operar seus próprios modelos em infraestrutura local. Essa mesma lógica se estende a organizações privadas que lidam com dados sensíveis e optam por soluções \textit{self-hosted} para manter controle sobre suas informações. Nesse ecossistema de implantação local, os mecanismos de segurança não vêm embutidos em um provedor externo: é responsabilidade de quem opera o modelo garantir que ele não produza ou processe conteúdos nocivos. Para isso, a comunidade de pesquisa desenvolveu os \textit{guardrails}, camadas de proteção que inspecionam as entradas e saídas de um LLM e sinalizam ou bloqueiam conteúdos potencialmente prejudiciais \cite{kumar2025freelunchguardrails}.

Em maio de 2026, o sistema Galileu, ferramenta de inteligência artificial desenvolvida pelo Tribunal Regional do Trabalho da 4\textordfeminine{} Região e nacionalizada pelo Conselho Superior da Justiça do Trabalho, identificou uma tentativa de manipulação em uma petição inicial processada pela 3\textordfeminine{} Vara do Trabalho de Parauapebas \cite{trt4galileu2026}. O documento continha trechos ocultos com instruções direcionadas ao próprio sistema, orientando-o a contestar a petição de forma superficial e a não impugnar os documentos apresentados, caracterizando um ataque de \textit{prompt injection}. Embora esse tipo de ataque esteja fora do escopo deste trabalho, o episódio evidencia que sistemas baseados em LLMs já estão implantados em contextos de alta criticidade no Brasil, operando sobre documentos redigidos em português, e que tentativas de manipulação desses sistemas já ocorrem em contextos reais. O episódio expõe uma fragilidade mais ampla: se sistemas de IA já enfrentam tentativas de manipulação em contextos reais no Brasil, em português, é igualmente necessário questionar se as próprias camadas de proteção desses sistemas são robustas nesse idioma. Uma questão fundamental, portanto, ainda não foi respondida: os \textit{guardrails} utilizados para proteger LLMs em produção classificam conteúdo prejudicial de forma confiável quando o texto está em português brasileiro? E o que acontece com essa confiabilidade quando esses modelos são comprimidos para se adequar às restrições de infraestrutura do mundo real?
	
A eficácia dos \textit{guardrails} em cenários multilíngues, no entanto, é uma preocupação recente e ainda pouco explorada. A literatura demonstra que LLMs treinados predominantemente em inglês tendem a apresentar vulnerabilidades de segurança quando expostos a entradas em outros idiomas: \citeonline{yong2024lowresource} mostram empiricamente que traduzir prompts nocivos para idiomas com menos recursos de treinamento é um método eficaz de contornar os mecanismos de segurança de modelos como o GPT-4, enquanto \citeonline{deng2024multilingual} documentam que a eficácia de filtros de toxicidade sofre degradação significativa quando os ataques são realizados em idiomas distintos do inglês. No contexto específico dos \textit{guardrails}, \citeonline{yang2024benchmarking} conduziram o estudo mais abrangente até o momento, avaliando \textit{guardrails} em mais de dez idiomas e concluindo que essas ferramentas ainda são ineficazes no tratamento de toxicidade multilíngue. Ainda assim, o português brasileiro não figura entre os idiomas avaliados em nenhum desses trabalhos. Recentemente, \citeonline{sousa-etal-2026-curupira} deram um primeiro passo ao desenvolver um guardrail específico para PT-BR, evidenciando que a comunidade começa a reconhecer essa lacuna, mas o desempenho desses modelos em português brasileiro, tanto em sua forma original quanto sob efeito de quantização, permanece sem investigação sistemática.

Para viabilizar a implantação de LLMs em infraestrutura local, uma prática amplamente adotada é a quantização. A técnica reduz a precisão numérica dos pesos do modelo, de 16 para 8 ou 4 bits, por exemplo, diminuindo o consumo de memória e aumentando o \textit{throughput} do sistema \cite{lin2024awq, frantar2023gptq}. Essa compressão é especialmente relevante no contexto dos \textit{guardrails}, que por operarem em paralelo ao modelo principal no mesmo pipeline de inferência, impõem restrições adicionais de memória e latência ao sistema \cite{kumar2025freelunchguardrails}. Contudo, a quantização não é gratuita: \citeonline{kharinaev2025quantization} demonstram que métodos de quantização pós-treinamento podem degradar o alinhamento de segurança de LLMs, com variações significativas entre métodos e níveis de precisão, e sem que um único método se destaque consistentemente. Esse achado levanta uma preocupação direta para o uso de \textit{guardrails} quantizados em cenários multilíngues. O efeito da compressão sobre modelos especializados em classificação de conteúdo prejudicial, especialmente quando operando em português brasileiro, permanece sem investigação sistemática na literatura, configurando a lacuna central que este trabalho se propõe a preencher.